\chapter{Atomic Data Sources}
\label{app:data}

Chapter~\ref{ch:atoms} explains what each rate in $\Lmat$ represents and why it
matters. This appendix records where each one came from, which script converted
it, and what it is stored in, so that any coefficient in the matrix can be
traced to a published source without reading the code. It also states which
coefficients are \emph{not} independent measurements, because that is the single
most important thing to know when reading a validation table: a check between
two quantities derived from one another cannot fail.

\section{The stack, end to end}

\begin{table}[htbp]
  \centering
  \caption[Provenance of every rate in the matrix]{Every coefficient entering
    $\Lmat$ or $\Svec$, with its source, the script that produced it, and the
    file it is stored in. Paths are relative to \texttt{data/processed/}.}
  \label{tab:data_provenance}
  \small
  \begin{tabular}{@{}p{2.6cm}p{3.5cm}p{3.2cm}p{4.3cm}@{}}
    \toprule
    quantity & source & script & stored in \\
    \midrule
    Excitation, \newline $n_{\rm upper} \le 10$
      & Convergent close-coupling cross sections
        \cite{BrayStelbovics1992,Bray2002}
      & \texttt{compute\_K\_CCC.py}
      & \texttt{collisions/ccc/} \newline \texttt{K\_CCC\_exc\_table.npy} \\
    \addlinespace
    Excitation, \newline $n = 11$--$15$
      & Vriens \& Smeets semi-empirical formula
        \cite{VriensSmeets1980}, in the form implemented by
        \cite{Hartgers2001}
      & \texttt{compute\_K\_VS.py}
      & \texttt{collisions/vs/} \newline \texttt{K\_VS\_exc\_table.npy} \\
    \addlinespace
    De-excitation
      & \emph{not independent}: detailed balance applied to the excitation
        rates above
      & \texttt{compute\_K\_CCC.py}, \newline \texttt{compute\_K\_VS.py}
      & \texttt{K\_deexc\_full.npy} \\
    \addlinespace
    Ionisation
      & CCC total ionisation cross sections, Maxwell-averaged, for
        $n \le 9$; Lotz's formula \cite{Lotz1968} for the bundled
        $n = 10$--$15$, which overestimates CCC by a factor $4$--$8$
      & \texttt{compute\_K\_TICS.py}, \newline \texttt{ionization\_rates.py}
      & \texttt{collisions/tics/} \newline \texttt{K\_ion\_final.npy} \\
    \addlinespace
    Radiative \newline recombination
      & Johnson's three-term Gaunt-factor formula \cite{Johnson1972},
        after \cite{Seaton1959}
      & \texttt{recombination\_rates.py}
      & \texttt{recombination/} \newline \texttt{alpha\_RR\_resolved.npy} \\
    \addlinespace
    Three-body \newline recombination
      & \emph{not independent}: detailed balance applied to the ionisation
        rates above
      & \texttt{recombination\_rates.py}
      & \texttt{recombination/} \newline \texttt{alpha\_3BR\_resolved.npy} \\
    \addlinespace
    Einstein $A$ \newline coefficients
      & Hoang-Binh's exact hydrogenic radial integrals \cite{HoangBinh1990},
        distributed as \textsc{aduu} \cite{HoangBinh2005}. The producing script
        names a 1993 paper, for which this project holds no reference; the
        discrepancy is unresolved
      & \texttt{radiative\_rates.py}
      & \texttt{Radiative/} \newline \texttt{A\_resolved.npy} \\
    \addlinespace
    Proton-impact \newline $\ell$-mixing
      & PSM20 Debye-cutoff formula \cite{Badnell2021}
      & \texttt{compute\_lmix.py}
      & \texttt{lmix/K\_lmix.npy} \\
    \addlinespace
    Assembled \newline operator
      & the eight rows above
      & \texttt{assemble\_cr\_matrix.py}
      & \texttt{cr\_matrix/L\_grid.npy}, \newline \texttt{S\_grid.npy},
        \texttt{L\_meta.csv} \\
    \addlinespace
    \midrule
    ADAS SCD96, \newline ACD96
      & \emph{comparison only; does not enter $\Lmat$}
        \cite{Summers2006,ADAS}
      & \texttt{parser\_adasf11.py}, \newline \texttt{prepare\_adas.py}
      & \texttt{adas/} \newline \texttt{SCD96\_interpolated.csv} \\
    \bottomrule
  \end{tabular}
\end{table}

\section{Which checks can fail, and which cannot}

Two rows of Table~\ref{tab:data_provenance} are marked \emph{not independent},
and the consequence is worth stating plainly rather than leaving for a reader to
infer. De-excitation is obtained from excitation by detailed balance, and
three-body recombination from ionisation by detailed balance. A numerical test
that those relations hold therefore tests the arithmetic that imposed them and
nothing else. Chapter~\ref{ch:validation} grades Gate~A \emph{tautological} for
exactly this reason: it passes at better than \SI{0.01}{\percent}, and it would
pass on corrupted input.

The checks that can fail are those between quantities with independent origins:
the excitation rates against R-matrix-with-pseudostates calculations, the raw
cross sections against microscopic reversibility, the assembled operator's
effective coefficients against ADAS, and the population coefficients against
Fujimoto's tabulation. The first two are graded \emph{severe} in
Table~\ref{tab:ladder} and the last two are among its four external
comparisons; of those four, one fails and one is an open disagreement.

\section{The seam at \texorpdfstring{$n = 10$}{n = 10}}

The excitation data change character between $n_{\rm upper} = 10$ and
$n_{\rm upper} = 11$: quantum close-coupling below, a semi-empirical formula
above. The seam is a property of the available data and not a modelling choice.
Its consequence for this thesis is confined to the bundled shells, which the
observable does not read: $\ffrac{3}-\ffrac{4}$ is built from $n = 3$ and
$n = 4$, seven and six shells below the seam.

\section{Two constants, and why both appear}

The hydrogen ionisation energy enters the pipeline twice, with two different
values, and this is deliberate. Collisional thresholds use the Rydberg energy
$R_\infty hc = \SI{13.605693122990}{\electronvolt}$, because that is the energy
scale on which the CCC and TICS cross sections are tabulated. Radiative
transition energies use the measured ionisation energy of the hydrogen atom,
$\SI{13.598434599702}{\electronvolt}$, which includes the reduced-mass and
relativistic corrections. Both are recorded in \texttt{L\_meta.csv} as
\texttt{IH\_collision\_eV} and \texttt{IH\_radiative\_eV}. They differ by
\SI{0.053}{\percent}. Appendix~\ref{app:states} states the hazard this creates for
detailed-balance tests, and Chapter~\ref{ch:atoms} the physics of why both are
correct in their own place.
